% Deep Learning chapter (M). Compile from report/ (tables in report/tables/, written by
% `python -m experiments.report_tables`; E5 tables by experiments/e5_accuracy.py).
% Rule 11: every number below is \input from a generated table or copied from a results file
% named in the comment beside it. Rule 10: cite only CLAUDE.md §16.1 papers and docs/refs files.
% Bibliography keys used (add to the shared .bib): gibbs2021aci, tu2026, bastani2022mvp,
% tibshirani2019, fannjiang2022, choi2026, li2021fno, pdebench2022, deepoheatv2,
% zrnic2024, angelopoulos2021.

\chapter{A calibrated neural surrogate gate}
\label{ch:dl}

This chapter covers the learning half of TwinHeat: a neural surrogate that predicts the steady
temperature of a four-layer board in milliseconds (Section~\ref{sec:dl-surrogate}), and a gate
that decides, design by design, whether that prediction may be trusted or must be checked by the
exact solver (Section~\ref{sec:dl-gate}). The gate's statistical guarantee is stated in
Section~\ref{sec:dl-guarantee} and tested in Section~\ref{sec:dl-experiments}.

\section{Problem and data}
\label{sec:dl-data}
A design is a $100\times100$~mm board with four copper layers, up to four chips on the top
layer, thermal vias, convection on both faces and optional heat-sunk edges. The quantity of
interest is the temperature rise $\theta$ above ambient on a $256\times256\times4$ grid
(262{,}144 unknowns), given by the finite-volume model of Chapter~\ref{ch:pagp}. The design
task used by the gate experiments varies 36 numbers: four chip positions, six copper rectangles
on layer~1 and the via density under each chip; everything else is fixed
(\texttt{config/design\_task.yaml}, decision D2).

Dataset v1 (\texttt{docs/decisions.md}, ``Dataset v1 split seeds'') has 6{,}000 training,
1{,}000 validation, 1{,}000 calibration and 1{,}000 test boards, half of the training boards
drawn from the 36-dimensional design space, plus out-of-distribution (OOD) sets of 500 boards
that each move one factor out of the training range (chip count, power, copper shapes, thin
traces, 3~oz copper, convection, boundary conditions, dense vias, chips on the bottom face).
Every field is solved to a relative residual of $10^{-10}$ and stored with SHA-256 hashes;
loaders refuse a hash mismatch, and the test and OOD splits are read only by the evaluation
scripts.

\section{The surrogate}
\label{sec:dl-surrogate}
The main model is a Fourier neural operator (FNO) \cite{li2021fno} that maps 13 input channels
(per-layer sheet conductance, chip power density, via density per gap, convection coefficients,
the heat-sunk-edge mask and two coordinate channels) to the four layer temperatures. It has
4~Fourier layers of width~48 with 24 modes. Five members, differing in seed and data order,
form an ensemble whose mean $\mu$ is the prediction and whose spread $\sigma$ is one of the
uncertainty signals. Each member is pre-trained at $128^2$ for 100 epochs and fine-tuned at
$256^2$ for 50 epochs, with a relative-$L_2$ loss plus 0.1 times the error on the 64 hottest
cells, and early stopping on the validation hotspot error (\texttt{docs/decisions.md},
``Ensemble training settings''). A U-Net and an FNO with a local $3\times3$ branch are trained
the same way as baselines. The FNO implementation was first checked against the published
Darcy-flow result on PDEBench \cite{pdebench2022}.

Table~\ref{tab:e5} gives the test accuracy (E5). The ensemble's mean relative $L_2$ error is
5.41\,\% % experiments/results/e5/go_no_go.md
against the 3\,\% target of the week-5 go/no-go, so that criterion was missed; its hotspot error,
0.764~K, % experiments/results/e5/go_no_go.md
meets the 1.5~K target. The team decided to continue with the ensemble at $256^2$ without
retraining (\texttt{docs/decisions.md}, week-5 go/no-go), because the gate acts on the peak
temperature, where the error is small, and the guarantee does not depend on model accuracy, only
the cost does. A single FNO with the local branch was more accurate on test than the ensemble;
the ensemble was kept for its spread. Table~\ref{tab:e5-ood} gives the OOD results: the error
grows on the sets that move power, convection and boundaries, and the feature-space OOD score
flags those sets far more often than the test set. The bottom-chip set is the one failure the
score does not catch: its error is very large and its flag rate is not.

\begin{table}[t]\centering\small
% E5 (experiments/e5_accuracy.py); every number is computed by the script.
\begin{tabular}{lrrrrr}
\hline
Model & rel.\ $L_2$ mean & rel.\ $L_2$ P95 & peak err.\ (K) & hotspot err.\ (K) & peak loc.\ (mm) \\
\hline
ensemble & 0.0541 & 0.1307 & 0.799 & 0.764 & 1.51 \\
fno\_pre128\_at\_256 & 0.2767 & 0.9555 & 2.029 & 1.858 & 3.55 \\
unet & 0.0543 & 0.1007 & 0.767 & 0.710 & 2.12 \\
fno\_local & 0.0478 & 0.1079 & 0.651 & 0.612 & 1.81 \\
member1 & 0.0640 & 0.1533 & 0.872 & 0.823 & 2.42 \\
member2 & 0.0603 & 0.1434 & 0.857 & 0.819 & 2.09 \\
member3 & 0.0618 & 0.1419 & 0.910 & 0.850 & 2.14 \\
member4 & 0.0610 & 0.1415 & 0.888 & 0.864 & 2.01 \\
member5 & 0.0625 & 0.1456 & 0.857 & 0.814 & 2.13 \\
\hline
\end{tabular}

\begin{tabular}{lrrrr}
\hline
Split (ensemble) & rel.\ $L_2$ mean & hotspot err.\ (K) & peak loc.\ (mm) & OOD flag rate \\
\hline
test & 0.0541 & 0.764 & 1.51 & 0.034 \\
ood\_1 & 0.0536 & 1.781 & 2.73 & 0.206 \\
ood\_2 & 0.1160 & 5.502 & 4.40 & 0.858 \\
ood\_3 & 0.0508 & 0.723 & 1.81 & 0.038 \\
ood\_4 & 0.0499 & 0.797 & 1.69 & 0.040 \\
ood\_6 & 0.1166 & 0.980 & 2.18 & 0.518 \\
ood\_7 & 0.1585 & 4.273 & 3.06 & 0.726 \\
ood\_3oz & 0.0681 & 1.126 & 2.03 & 0.050 \\
ood\_dense\_vias & 0.0616 & 0.888 & 1.90 & 0.046 \\
ood\_bottom\_chips & 0.9374 & 24.312 & 42.45 & 0.002 \\
\hline
\end{tabular}

\caption{E5 surrogate accuracy on the test split and the OOD sets (tag of the E5 run:
\texttt{experiments/results/e5/run\_info.json}). Top: models; bottom: the ensemble per split, with
the rate at which the OOD score flags a board.}
\label{tab:e5}\label{tab:e5-ood}
\end{table}

\section{The gate}
\label{sec:dl-gate}
For a proposed design the server computes the ensemble prediction and a label-free scale
\[
  s = r_{\mathrm{hot}} + \lambda\,\sigma_{\mathrm{hot}} + \varepsilon ,
\]
where $r_{\mathrm{hot}}$ is the mean absolute residual $|(A\mu-b)/\operatorname{diag}A|$ of the
discrete heat equation over the 64 hottest cells of $\mu$ (one stencil pass, in kelvin) and
$\sigma_{\mathrm{hot}}$ the ensemble spread there. The gate accepts the prediction when
$B = q_t\,s \le \tau$ and otherwise flags the design to the exact solver. Everything the gates
differ in is how $q_t$ is set:
\begin{itemize}
\item G0: no gate; G1: the DeepOHeat-v2 rule \cite{deepoheatv2}, flagging when $r_{\mathrm{hot}}$
exceeds a moving percentile of its last 100 values, with the flag rate matched to G5's
exact-solve budget;
\item G2, G3: split conformal and weighted split conformal \cite{tibshirani2019,fannjiang2022,choi2026},
calibrated once on 1{,}000 labelled designs drawn from the design-space density;
\item G4: adaptive conformal inference (ACI) \cite{gibbs2021aci} updated with the labels of
flagged designs only, an ablation;
\item G5: ACI updated only on randomly audited designs; G6: the same audits feeding one threshold
per group through the MultiValid predictor \cite{bastani2022mvp}; G7: audits whose probability
grows with the scale, with a floor $p_{\min}$ \cite{zrnic2024}.
\end{itemize}
An audit is a coin $A_t\sim\mathrm{Bernoulli}(p)$ drawn by the server \emph{after} the design is
committed and never shown to the proposer (invariant B8); an audited design is solved exactly and
its error updates $q_t$. Designs come from three automated proposers with three aggressiveness
levels each: a tilted MCMC sampler that favours low predicted peaks (P1), simulated annealing
(P2) and batched gradient descent through the surrogate (P3). A logged-human proposer (P4) was
planned but needs the user study, which was not run.

\section{Guarantee}
\label{sec:dl-guarantee}
Let $W_T$ be the fraction of $T$ proposals that were accepted while their true peak error exceeded
$\tau$. If the proposer's design does not depend on the current coin and $q_t$ is updated only on
audited steps by the ACI rule with step $\gamma$ and level $\alpha$, then with probability at
least $1-\delta$
\[
  W_T \;\le\; \alpha + \frac{C}{\gamma p T} + \frac{2}{p}\sqrt{\frac{\ln(4/\delta)}{T}},
  \qquad C = \max(\alpha_1, 1-\alpha_1) + \gamma ,
\]
for any proposer, adversarial or human (Lemma~1, Appendix~\ref{app:lemmas}). The deterministic
control of the audited steps is Gibbs and Cand\`es' Proposition~4.1 \cite{gibbs2021aci}, applied
to the revealed subsequence as Tu and Giesecke do \cite{tu2026}; the transfer to all steps in
expectation is Tu and Giesecke's Theorem~5.9. Our contributions are the first step (an accepted
large error is a miscoverage) and the high-probability transfer, which uses the fact that the
miscoverage indicator is fixed before the coin is drawn. When an audit stops early with a
certified interval $[T^-,T^+]$ for the true peak (Chapter~\ref{ch:pagp}), the conservative score
$S^+ = \max(|\hat T - T^-|, |\hat T - T^+|)/s$ keeps the bound (Lemma~1$'$).

The bound is a long-run average, not a per-step statement, and it is loose at our run lengths:
at $T = 2{,}000$, $p = 0.1$ and $\gamma = 0.05$ it exceeds~1, so it cannot by itself separate
the gates. % report/lemma1_check_n31.md section 2; docs/decisions.md D4
We therefore report the measured rate $M_1 = W_T$ against $\alpha$ directly, and the bound beside
it.

\section{Experiments}
\label{sec:dl-experiments}
\paragraph{Setup.} 45 proposer runs (P1--P3, three levels, five seeds) produced 92{,}461
proposals; % docs/decisions.md (D4 amendment): the count on rvcse
with the 1{,}000 calibration designs these contain 92{,}311 distinct designs,
% experiments/oracle_report.md (93,461 input rows, 92,311 unique)
and every one was labelled with an exact solve to a relative residual of $10^{-8}$, cross-checked on 0.5\,\% of the designs by a second solver (largest
difference $2.7\times10^{-8}$~K, \texttt{experiments/oracle\_report.md}). Every gate is then
replayed offline on the same trajectories with the same coins. All hyperparameters were frozen
before the runs ($\alpha = 0.1$, $\tau = 1$~K, $\gamma = 0.05$, window 200, headline $p = 0.1$;
\texttt{experiments/configs/frozen.yaml}, decision D4), and every number comes from a tagged
commit.

\paragraph{The pre-registered threshold was too loose to test the gates.} At $\tau = 1$~K the
ensemble's peak error exceeded $\tau$ on only 0.1\,\% of the proposals (median error 0.115~K),
% docs/decisions.md, D4 amendment
so even with no gate the accepted-large-error rate is far below $\alpha$
(Table~\ref{tab:e7-tau1}). Every gate passes trivially and the experiment cannot rank them. This
is reported as the pre-registered result.

% generated by experiments/report_tables.py from experiments/results/e7/e7_runs.csv
\begin{table}[t]\centering\small
\begin{adjustbox}{max width=\linewidth}\begin{tabular}{llrrrr}
\hline
Gate & & $p$ & mean $M_1$ & runs with $M_1>\alpha$ & exact solves / proposal \\
\hline
G0 & no gate & -- & 0.0011 & 0/45 & 0.001 \\
G1 & residual percentile (DeepOHeat-v2) & -- & 0.0005 & 0/45 & 0.105 \\
G2 & split conformal & -- & 0.0005 & 0/45 & 0.012 \\
G3 & weighted conformal & -- & 0.0005 & 0/45 & 0.164 \\
G4 & ACI, flagged labels only & -- & 0.0006 & 0/45 & 0.008 \\
G5 & audited ACI & 0.05 & 0.0005 & 0/45 & 0.065 \\
G5 & audited ACI & 0.1 & 0.0005 & 0/45 & 0.117 \\
G5 & audited ACI & 0.2 & 0.0005 & 0/45 & 0.221 \\
G6 & audited multivalid & -- & 0.0008 & 0/45 & 0.102 \\
G7 & audited ACI, risk-adaptive & -- & 0.0005 & 0/45 & 0.115 \\
\hline
\end{tabular}\end{adjustbox}
\caption{E7 at the pre-registered $\tau = 1$~K (tag v3-freeze). $\alpha = 0.1$; 450 runs in all, 45 per row (P1--P3 $\times$ 3 aggressiveness levels $\times$ 5 seeds).}
\label{tab:e7-tau1}
\end{table}


\paragraph{Sensitivity run at $\tau = 0.2$~K.} We added one run, decided after seeing the
error distribution and logged before running (D4 amendment), at $\tau = 0.2$~K, where 25\,\% of
the proposals are large errors, so that the gates must work to stay below $\alpha = 0.1$
(Table~\ref{tab:e7-tau02}). Without a gate every run exceeds $\alpha$. The audited ACI gate G5
never exceeds $\alpha$ in any of its 135 runs, at all three audit rates, and neither does G7.
ACI fed only with the labels of flagged designs (G4) exceeds $\alpha$ in 7 of 45 runs: it
calibrates on a selected population, the bias Tu and Giesecke quantify (Proposition~B.14 of
\cite{tu2026}), which is why only audits update the threshold. Weighted conformal (G3) exceeds
$\alpha$ in 4 runs, all with the tilted sampler P1, for which its weights are exact; we do not
have an explanation and report it as observed.

% generated by experiments/report_tables.py from experiments/results/e7_tau0.2/e7_runs.csv
\begin{table}[t]\centering\small
\begin{tabular}{llrrrr}
\hline
Gate & & $p$ & mean $M_1$ & runs with $M_1>\alpha$ & exact solves / proposal \\
\hline
G0 & no gate & -- & 0.2581 & 45/45 & 0.001 \\
G1 & residual percentile (DeepOHeat-v2) & -- & 0.0403 & 0/45 & 0.832 \\
G2 & split conformal & -- & 0.0136 & 0/45 & 0.927 \\
G3 & weighted conformal & -- & 0.0314 & 4/45 & 0.883 \\
G4 & ACI, flagged labels only & -- & 0.0625 & 7/45 & 0.670 \\
G5 & audited ACI & 0.05 & 0.0177 & 0/45 & 0.907 \\
G5 & audited ACI & 0.1 & 0.0201 & 0/45 & 0.900 \\
G5 & audited ACI & 0.2 & 0.0230 & 0/45 & 0.896 \\
G6 & audited multivalid & -- & 0.0903 & 16/45 & 0.669 \\
G7 & audited ACI, risk-adaptive & -- & 0.0201 & 0/45 & 0.896 \\
\hline
\end{tabular}
\caption{E7 at the sensitivity $\tau = 0.2$~K (tag v3.1-freeze, D4 amendment). $\alpha = 0.1$; 450 runs in all, 45 per row (P1--P3 $\times$ 3 aggressiveness levels $\times$ 5 seeds).}
\label{tab:e7-tau02}
\end{table}


The price is the exact-solve budget. At $\tau = 0.2$~K the scale $s$ is not sharp enough to
accept most designs, so G5 sends about 90\,\% of them to the exact solver
(Table~\ref{tab:e7-tau02}, last column). The residual-percentile gate G1, run at the same budget,
also stays below $\alpha$, so the hypothesis that G1 fails as the proposers become more aggressive
(H1) is not supported: at a matched budget this large, both gates are safe. A sharper
label-free scale would lower the budget; the guarantee does not depend on it.

\paragraph{Per-group validity (E15, H8).} Table~\ref{tab:e15-tau02} splits the accepted-error
rate by proposer, aggressiveness, OOD-score tercile and hotspot position. G5 stays below
$\alpha$ in all but 2 of its 277 group-runs, both in small groups
% counts from report/tables/e15_tau02.tex (G5 column)
(\texttt{edge=1}, \texttt{ood=3}). The
multivalid gate G6, designed for per-group control, exceeds $\alpha$ in many groups, most often
with simulated annealing (P2). H8 is not supported; G6's per-group guarantee was a target for a
proof that was not completed and is reported as empirical. The likely causes are the small
number of audits per group and bucket at $p = 0.1$ and the untuned MultiValid settings
(D4, D9).

% generated by experiments/report_tables.py from experiments/results/e15_tau0.2/e15_groups.csv
\begin{table}[t]\centering\small
\begin{adjustbox}{max width=\linewidth}\begin{tabular}{lrr}
\hline
Group & G5: runs with $M_{1,g}>\alpha$ & G6: runs with $M_{1,g}>\alpha$ \\
\hline
aggr=1 & 0/15 & 5/15 \\
aggr=2 & 0/15 & 5/15 \\
aggr=3 & 0/15 & 6/15 \\
edge=0 & 0/45 & 16/45 \\
edge=1 & 1/32 & 5/32 \\
ood=1 & 0/45 & 14/45 \\
ood=2 & 0/45 & 22/45 \\
ood=3 & 1/20 & 14/20 \\
proposer=P1 & 0/15 & 1/15 \\
proposer=P2 & 0/15 & 15/15 \\
proposer=P3 & 0/15 & 0/15 \\
\hline
\end{tabular}\end{adjustbox}
\caption{E15 per-group accepted-error rate at $\tau = 0.2$~K, $p = 0.1$. $\alpha = 0.1$.}
\label{tab:e15-tau02}
\end{table}


\paragraph{Risk-adaptive audits (E16, H9).} Auditing high-uncertainty designs more often (G7)
keeps the rate below $\alpha$ in every run but saves at most about 10\,\% of the audits relative
to uniform audits (Table~\ref{tab:e16-tau02}), short of the 30\,\% that H9 asked for. Uniform
audits remain the default. G7's update has no cited validity argument and is reported as an
ablation only.

% generated by experiments/report_tables.py from experiments/results/e16_tau0.2/e16_runs.csv
\begin{table}[t]\centering\small
\begin{tabular}{lrrrrr}
\hline
Gate & $\bar p$ & $p_{\min}$ & audits / run & mean $M_1$ & runs with $M_1>\alpha$ \\
\hline
G5 & 0.05 & -- & 101.5 & 0.0177 & 0/45 \\
G5 & 0.1 & -- & 206.0 & 0.0201 & 0/45 \\
G5 & 0.2 & -- & 418.8 & 0.0230 & 0/45 \\
G7 & 0.05 & 0.02 & 92.4 & 0.0186 & 0/45 \\
G7 & 0.05 & 0.05 & 104.4 & 0.0182 & 0/45 \\
G7 & 0.1 & 0.02 & 187.2 & 0.0201 & 0/45 \\
G7 & 0.1 & 0.05 & 187.2 & 0.0201 & 0/45 \\
G7 & 0.2 & 0.02 & 377.2 & 0.0224 & 0/45 \\
G7 & 0.2 & 0.05 & 377.2 & 0.0224 & 0/45 \\
\hline
\end{tabular}
\caption{E16 risk-adaptive audits (G7) against uniform audits (G5) at $\tau = 0.2$~K. $\alpha = 0.1$.}
\label{tab:e16-tau02}
\end{table}


\paragraph{Certified audits (E14, H7).} An audit need not solve to $10^{-8}$: it can stop
as soon as its certified score interval lies on one side of the threshold, or is narrower than
$\varepsilon_S$ (Lemma~1$'$; Chapter~\ref{ch:pagp}). Replaying G5 with such audits, from
certified traces of all 9{,}176 audited designs % experiments/audit_traces.py run, tag v3.4-freeze
computed at tag \texttt{v3.4-freeze} with every interval checked against the oracle, uses about
48\,\% fewer audit iterations than exact audits at every $\varepsilon_S$ of the sweep, and no run
exceeds $\alpha$ (Table~\ref{tab:e14-tau02}). The saving at the frozen $\varepsilon_S = 0.1$ is
47.7\,\%, % experiments/results/e14_tau0.2/e14_verdicts.json
just under the 50\,\% that H7 asked for, so H7 is not supported. Because $S^+\ge S$, certified
audits make the threshold slightly conservative, which lowers the accepted-error rate further
than exact audits do. The audited designs and the stopping rule do not depend on $\tau$, so the
same saving holds at $\tau = 1$~K (\texttt{experiments/results/e14}).

% generated by experiments/report_tables.py from experiments/results/e14_tau0.2/e14_runs.csv
\begin{table}[t]\centering\small
\begin{tabular}{lrrrrr}
\hline
Audits & $\varepsilon_S$ & audit iterations & saved vs.\ $10^{-8}$ & mean $M_1$ & runs with $M_1>\alpha$ \\
\hline
certified & 0.05 & 6{,}352{,}580 & 47.8\% & 0.0086 & 0/45 \\
certified & 0.1 & 6{,}365{,}930 & 47.7\% & 0.0082 & 0/45 \\
certified & 0.2 & 6{,}304{,}660 & 48.2\% & 0.0072 & 0/45 \\
exact & -- & 12{,}161{,}420 & -- & 0.0201 & 0/45 \\
\hline
\end{tabular}
\caption{E14 certified audits (Lemma 1$'$) against audits solved to $10^{-8}$, at $\tau = 0.2$~K (tag v3.4-freeze). $\alpha = 0.1$; audit iterations summed over the 45 runs of G5 at $p = 0.1$.}
\label{tab:e14-tau02}
\end{table}


\section{Limitations}
\begin{itemize}
\item The pre-registered $\tau = 1$~K was too loose for the experiment to rank the gates; the
informative run at $\tau = 0.2$~K was chosen after seeing the error distribution. Both are
reported.
\item The Lemma~1 bound is vacuous at our run lengths and audit rates; the evidence is the measured
rate against $\alpha$ over 45 runs per gate, which is empirical.
\item At $\tau = 0.2$~K the gates are safe but expensive (about 90\,\% exact solves); the
efficiency comparison at a fixed budget with feedback to the proposer (E8, H2) was not run.
\item The surrogate missed its relative-$L_2$ target, and the OOD score does not flag the
bottom-chip set.
\item No human proposer (P4) and no user study (E12, H4): the pilot needs ethics approval. The
Darcy replication (E9, H3) was cut.
\item G6's per-group bound is unproved and its empirical per-group control failed here; G3's
failures with P1 are unexplained.
\end{itemize}
